\documentclass[12pt,a4paper]{article}

% Load Configurations
% Packages Configuration
\usepackage[utf8]{inputenc}
\usepackage[T1]{fontenc}
\usepackage{mathptmx} % Times New Roman Font

% Indonesian Language Support
\usepackage[indonesian]{babel}

% Layout & Geometry
\usepackage[margin=1in]{geometry}
\usepackage{microtype}
\usepackage{setspace}

% Math & Symbols
\usepackage{amsmath}
\usepackage{amssymb}
\usepackage{amsfonts}

% Figures & Tables
\usepackage{graphicx}
\usepackage{booktabs}
\usepackage{tabularx}
\usepackage{array}
\usepackage{longtable}

% Bibliography
\usepackage[numbers,sort&compress]{natbib}

% Table of Contents Formatting
\usepackage{tocloft}

% Section Title Formatting
\usepackage{titlesec}

% Hyperlinks (All links black & hidelinks)
\usepackage{hyperref}

% Formatting Configuration

% Line Spacing (1.5 Spacing)
\onehalfspacing

% Hyperref Setup (All Black / Hide Links)
\hypersetup{
    colorlinks=false,
    hidelinks,
    pdftitle={[Judul Penelitian Skripsi]},
    pdfauthor={[Nama Mahasiswa]}
}

% Section Titles: Centered, Bold, Uppercase, Unnumbered Label
% Subsections: Numbered (1.1, 1.2, 2.1, 2.2...)
\setcounter{secnumdepth}{2}

% Make every \section start on a new page
\newcommand{\sectionbreak}{\clearpage}

\titleformat{\section}[block]{\centering\Large\bfseries}{}{0pt}{}
\titleformat{\subsection}[block]{\large\bfseries}{\thesubsection\quad}{0pt}{}
\titleformat{\subsubsection}[block]{\bfseries}{\thesubsubsection\quad}{0pt}{}

% Table of Contents (Daftar Isi) Formatting
\renewcommand{\contentsname}{DAFTAR ISI}
\renewcommand{\cfttoctitlefont}{\hfill\Large\bfseries}
\renewcommand{\cftaftertoctitle}{\hfill\null}

% Dense dot leaders for all entries (titik-titik padat)
\renewcommand{\cftdotsep}{0.5}
\renewcommand{\cftsecleader}{\cftdotfill{\cftdotsep}}
\renewcommand{\cftsubsecleader}{\cftdotfill{\cftdotsep}}

% Remove section number for major section titles in TOC
\newcommand{\cftignore}[1]{}
\renewcommand{\cftsecpresnum}{\cftignore}
\setlength{\cftsecnumwidth}{0pt}

% Section font in TOC: Bold
\renewcommand{\cftsecfont}{\normalfont\bfseries}
\renewcommand{\cftsecpagefont}{\normalfont\bfseries}

% Rename Bibliography Header to DAFTAR PUSTAKA
\addto\captionsindonesian{
  \renewcommand{\refname}{DAFTAR PUSTAKA}
}

% Paragraph Spacing
\setlength{\parskip}{0.5em}
\setlength{\parindent}{0pt}

% Custom LaTeX Commands

% Helper command for marking items as not reported
\newcommand{\notreported}{\textit{Not reported}}
\newcommand{\notverified}{\textit{Not verified}}

% Custom command for paper citations / references shortcut
\newcommand{\paperref}[2]{Paper #1 (\cite{#2})}


\begin{document}

% ==================== COVER PAGE (SESUAI TEMPLATE POLINEMA) ====================
\begin{titlepage}
    \centering
    \vspace*{0.2cm}
    
    {\Large\bfseries [JUDUL PENELITIAN SKRIPSI]\par}
    
    \vspace{1.0cm}
    
    {\large\bfseries STUDI LITERATUR\par}
    \vspace{0.3cm}
    {\large Disusun untuk Memenuhi Salah Satu Prasyarat Bimbingan Skripsi\par}
    
    \vspace{1.0cm}
    
    {\large\bfseries Oleh:\par}
    \vspace{0.2cm}
    {\large\bfseries [Nama Mahasiswa] \qquad NIM. [Nomor Induk Mahasiswa]\par}
    
    \vspace{0.8cm}
    
    \includegraphics[width=3.8cm]{assets/polinema-logo.png}
    
    \vfill
    
    {\large\bfseries PROGRAM STUDI SARJANA TERAPAN TEKNIK INFORMATIKA\par}
    {\large\bfseries JURUSAN TEKNOLOGI INFORMASI\par}
    {\large\bfseries POLITEKNIK NEGERI MALANG\par}
    {\large\bfseries 2026\par}
\end{titlepage}

\cleardoublepage
\pagenumbering{roman}
\setcounter{page}{1}

% ==================== ABSTRAK ====================
\section*{ABSTRAK}
\addcontentsline{toc}{section}{ABSTRAK}
[Tuliskan uraian abstrak studi literatur atau proposal penelitian skripsi Anda di sini. Abstrak memuat konteks latar belakang permasalahan, tujuan penelitian, metodologi yang dikaji, kesenjangan penelitian yang ditemukan, serta kontribusi atau arah penelitian yang diusulkan.]

\vspace{0.5cm}
\textbf{Kata Kunci}: [Kata Kunci 1], [Kata Kunci 2], [Kata Kunci 3], [Kata Kunci 4], Studi Literatur.

\newpage

% ==================== DAFTAR ISI ====================
\tableofcontents
\newpage

% ==================== ISI UTAMA (ARABIC PAGE NUMBERING) ====================
\pagenumbering{arabic}
\setcounter{page}{1}

\section{PENDAHULUAN}
\label{sec:pendahuluan}

\subsection{Latar Belakang}
Sektor akuakultur air tawar di Indonesia memegang peranan strategis dalam ketahanan pangan nasional dan pemberdayaan ekonomi masyarakat, dengan ikan lele (\textit{Clarias sp.}) sebagai salah satu komoditas unggulan utama. Namun, efisiensi operasional budidaya lele intensif sangat dibatasi oleh tingginya porsi biaya pakan yang mencapai 60\% hingga 70\% dari keseluruhan biaya produksi~\cite{Li2020}. Hingga saat ini, sebagian besar pembudidaya masih mengandalkan metode pemberian pakan konvensional berbasis jadwal waktu tetap (\textit{fixed-time scheduling}) atau estimasi visual manual secara subjektif.

Metode konvensional tersebut memicu dua permasalahan kritis:
\begin{enumerate}
    \item \textbf{Pemberian Pakan Berlebih (\textit{Overfeeding})}: Pelet pakan yang tidak terkonsumsi akan tenggelam dan mengalami dekomposisi organik di dasar kolam. Proses pembusukan ini meningkatkan konsentrasi amonia terlarut ($NH_3$) dan menurunkan kadar oksigen terlarut (\textit{Dissolved Oxygen}), sehingga memicu mortalitas massal dan memperburuk rasio konversi pakan (\textit{Feed Conversion Ratio} / FCR)~\cite{Zhou2018}.
    \item \textbf{Kekurangan Pakan (\textit{Underfeeding})}: Ikan lele memiliki sifat biologis kanibalistik yang sangat agresif ketika mengalami defisit energi atau kelaparan (\textit{starvation-stress})~\cite{Marliza2021}. Keterlambatan pemberian pakan secara langsung memicu agresi sesama ikan, menyebabkan penurunan biomassa panen secara signifikan.
\end{enumerate}

Untuk mewujudkan otomasi pemberian pakan adaptif (\textit{demand-feeding system}), teknologi visi komputer (\textit{computer vision}) non-kontak menjadi solusi yang sangat potensial untuk memantau respons nafsu makan ikan secara langsung tanpa menimbulkan stres fisik~\cite{Al-Abri2025,Yang2025}. Kendati demikian, implementasi visi komputer pada kolam budidaya lele luar ruangan (\textit{outdoor}) menghadapi sejumlah kendala teknis ekstrem:
\begin{enumerate}
    \item \textbf{Keberadaan Aerator Kontinu}: Berbeda dengan tangki laboratorium tertutup, kolam lele komersial dilengkapi aerator mekanis (kincir atau gelembung udara) yang beroperasi secara kontinu selama 24 jam untuk menjaga pasokan oksigen. Turbulensi dan riak air konstan ini menyebabkan metode deteksi gerak konvensional mendeteksi pergerakan palsu secara terus-menerus. Dengan demikian, prinsip dasar yang harus ditegakkan adalah: \textit{pergerakan permukaan air kolam tidak sama dengan indikasi ikan lapar}.
    \item \textbf{Karakteristik Temporal Kemunculan Sub-Detik}: Air kolam lele bersifat keruh (berwarna hijau pekat atau cokelat akibat alga dan lumpur), sehingga tubuh ikan tidak terdeteksi saat berada di kedalaman air. Ikan lele hanya muncul ke permukaan air ketika terstimulasi lapar, menyambar pakan, dan segera kembali menyelam ke dasar kolam. Kemunculan ini berlangsung dalam jendela waktu transien yang sangat singkat, berkisar antara 3 hingga 6 frame pada kecepatan rekam 10 FPS (durasi total hanya 0,3 hingga 0,6 detik).
    \item \textbf{Keterbatasan Komputasi Perangkat Tertanam (\textit{Edge Computing})}: Sistem pemantau harus diimplementasikan langsung di tepi kolam menggunakan perangkat berdaya rendah seperti Raspberry Pi 4B/5 tanpa dukungan kartu grafis (GPU) diskrit. Model deep learning konvensional yang memerlukan miliaran operasi komputasi tidak dapat berjalan secara stabil pada frekuensi $\ge 10$ FPS.
    \item \textbf{Keterbatasan Bandwidth Komunikasi}: Jaringan komunikasi di tambak rakyat umumnya memiliki keterbatasan kuota dan stabilitas sinyal. Sistem dituntut untuk memproses komputasi secara lokal di \textit{edge} dan hanya mentransmisikan telemetri metrik numerik ringkas (indeks deviasi gerak dan status lapar) melalui protokol IoT seperti MQTT, bukan mentransmisikan aliran video mentah (\textit{raw video streaming}).
\end{enumerate}

Berdasarkan latar belakang tersebut, penelitian ini bertujuan mengembangkan metode deteksi aktivitas gerak ikan lele berbasis visi komputer yang ringan, adaptif terhadap derau aerator, dan mampu mengekstraksi event kemunculan sub-detik secara real-time pada perangkat Raspberry Pi.

\subsection{Perumusan Masalah}
Berdasarkan latar belakang dan batasan lingkungan yang telah diuraikan, rumusan masalah dalam penelitian ini dirumuskan sebagai berikut:
\begin{enumerate}
    \item Bagaimana memodelkan dan mengompensasi derau gerak permukaan air yang ditimbulkan oleh aerator kontinu pada kolam luar ruangan agar tidak memicu deteksi positif palsu (\textit{false positive})?
    \item Bagaimana merancang algoritma deteksi gerakan spasio-temporal yang memiliki kompleksitas komputasi rendah sehingga mampu mengeksekusi analisis citra secara real-time ($\ge 10$ FPS) pada perangkat Raspberry Pi untuk menangkap event kemunculan lele berdurasi 3--6 frame?
    \item Bagaimana memetakan deviasi gerak permukaan air menjadi indeks kuantitatif aktivitas makan dan menghasilkan keputusan status lapar (\textit{hungry} vs. \textit{satiated}) yang akurat?
\end{enumerate}

\subsection{Tujuan Penelitian}
Tujuan yang hendak dicapai dalam penelitian ini adalah:
\begin{enumerate}
    \item Mengembangkan mekanisme pemodelan \textit{baseline} dinamis untuk mengeliminasi komponen derau gerak aerator pada citra kolam luar ruangan.
    \item Merancang dan mengimplementasikan algoritma deteksi gerak hibrida berbasis varians temporal dan perbedaan citra yang ringan untuk mengekstraksi lonjakan aktivitas makan sub-detik pada Raspberry Pi.
    \item Menghasilkan metrik kuantitatif indeks aktivitas makan ikan lele yang dapat ditransmisikan secara efisien melalui jaringan IoT sebagai acuan otomasi pakan.
\end{enumerate}

\subsection{Manfaat Penelitian}
Penelitian ini diharapkan memberikan kontribusi ganda:
\begin{itemize}
    \item \textbf{Manfaat Akademis}: Memberikan kontribusi teoritis dalam pemodelan latar belakang dinamis (\textit{dynamic background subtraction}) pada fluida bergolak dan ekstraksi event spasio-temporal transien berdaya rendah di bidang visi komputer terapan akuakultur.
    \item \textbf{Manfaat Praktis}: Menyediakan solusi teknologi tepat guna berbiaya rendah (\textit{low-cost edge AI}) bagi peternak lele lokal untuk mencegah pemborosan pakan, menjaga kualitas air kolam, dan menekan angka kanibalisme ikan.
\end{itemize}

\section{LANDASAN TEORI}
\label{sec:landasan-teori}

\subsection{Karakteristik Biologis dan Perilaku Makan Ikan Lele}
Ikan lele (\textit{Clarias sp.}) merupakan spesies ikan air tawar bertubuh memanjang dan tidak bersisik yang tergolong dalam ordo Siluriformes. Secara biologis, ikan lele memiliki organ pernapasan tambahan berupa \textit{arborescent organ} yang memungkinkannya bertahan pada perairan dengan konsentrasi oksigen terlarut rendah~\cite{Marliza2021}. Dalam kondisi lingkungan kolam intensif, lele memiliki karakteristik fotofobik (menghindari intensitas cahaya matahari terik) dan cenderung berdiam atau berenang perlahan di dasar perairan yang terlindung.

Ketika mengalami peningkatan nafsu makan atau kelaparan (\textit{starvation-stress}), pola perilaku renang lele mengalami transisi drastis dari kondisi diam (\textit{quiescent state}) menjadi fase pencarian makan aktif (\textit{active foraging})~\cite{Alahmad2025}. Sifat makan lele ditandai dengan gerakan meluncur vertikal berkecepatan tinggi menuju permukaan air, menyambar pakan dengan gerakan menyamping (\textit{surface thrashing}), memicu buih dan cipratan air, lalu menyelam kembali ke dasar kolam~\cite{Li2020}. Karakteristik kemunculan permukaan ini bersifat transien dan terlokalisasi, dengan rentang durasi tipikal hanya 3 hingga 6 frame pada laju tangkap kamera 10 FPS (0,3--0,6 detik).

\subsection{Metode Deteksi Gerakan dalam Visi Komputer}
Deteksi gerakan objek pada sekuens video merupakan proses identifikasi region spasial yang mengalami perubahan intensitas piksel signifikan seiring berjalannya waktu. Terdapat tiga paradigma klasik dan satu paradigma modern yang umum diterapkan dalam pemantauan visual akuakultur~\cite{Al-Abri2025,Yang2025}:

\subsubsection{Temporal Frame Differencing}
Metode perbedaan frame temporal (\textit{temporal frame differencing}) mengekstrak region gerak dengan menghitung selisih nilai intensitas absolut antara dua atau tiga frame yang berurutan secara berpasangan~\cite{Dong2024}. Untuk sekuens citra skala abu-abu $I(x,y,t)$, operasi selisih dua frame dinyatakan sebagai:
\begin{equation}
    D_2(x,y,t) = |I(x,y,t) - I(x,y,t-1)|
    \label{eq:frame-diff-2}
\end{equation}
Region gerak biner $M(x,y,t)$ diperoleh melalui penerapan ambang batas (\textit{thresholding}):
\begin{equation}
    M(x,y,t) = \begin{cases} 
        1, & \text{jika } D_2(x,y,t) > \tau \\
        0, & \text{lainnya}
    \end{cases}
    \label{eq:binary-diff}
\end{equation}
Metode ini memiliki keunggulan utama berupa kompleksitas komputasi yang sangat rendah berordo $O(N)$ (di mana $N$ adalah jumlah piksel), kebutuhan memori RAM yang minimal, serta sensitivitas instan tanpa jeda waktu akumulasi histori~\cite{Liu2014}. Namun, kelemahan utamanya adalah kerentanan ekstrem terhadap derau fluida bergerak kontinu seperti gelembung aerator.

\subsubsection{Optical Flow}
Aliran optik (\textit{optical flow}) memodelkan pola kecepatan semu dari pergerakan objek antara dua frame berurutan berdasarkan asumsi konsistensi kecerahan (\textit{brightness constancy assumption}):
\begin{equation}
    I(x,y,t) = I(x + \Delta x, y + \Delta y, t + \Delta t)
    \label{eq:brightness-constancy}
\end{equation}
Ekspansi deret Taylor orde pertama menghasilkan persamaan konstrain aliran optik dasar:
\begin{equation}
    \frac{\partial I}{\partial x} u + \frac{\partial I}{\partial y} v + \frac{\partial I}{\partial t} = 0 \quad \iff \quad I_x u + I_y v + I_t = 0
    \label{eq:optical-flow-constraint}
\end{equation}
di mana $(u,v) = \left(\frac{\Delta x}{\Delta t}, \frac{\Delta y}{\Delta t}\right)$ merepresentasikan komponen vektor kecepatan gerak horizontal dan vertikal. 

Pendekatan aliran optik terbagi menjadi dua kategori:
\begin{itemize}
    \item \textbf{Dense Optical Flow (Metode Gunnar Farnebäck)}: Menghitung vektor perpindahan pada seluruh piksel citra secara rapat~\cite{Fukae2023}. Pendekatan ini mampu menangkap dinamika kawanan ikan secara utuh, namun menuntut operasi matriks kuadrat terkecil lokal yang sangat intensif, menjadikannya tidak layak dieksekusi secara real-time pada prosesor berdaya rendah.
    \item \textbf{Sparse Optical Flow (Metode Lucas-Kanade)}: Hanya menghitung vektor gerak pada titik-titik sudut fitur terdefinisi (misalnya titik fitur Shi-Tomasi)~\cite{Zhang2024}. Pendekatan sparse jauh lebih ringan dan cocok diterapkan pada region deviasi gerak terpilih di perangkat edge.
\end{itemize}

\subsubsection{Adaptive Background Subtraction dan Gaussian Mixture Models}
Pengurangan latar belakang adaptif memisahkan objek latar depan (\textit{foreground}) dengan cara memodelkan dan memperbarui representasi latar belakang referensi $B(x,y,t)$ secara berkala~\cite{Aris2023}:
\begin{equation}
    B(x,y,t) = (1 - \alpha) B(x,y,t-1) + \alpha I(x,y,t)
    \label{eq:running-average}
\end{equation}
di mana $\alpha \in (0,1)$ adalah laju pembelajaran (\textit{learning rate}). Pada lingkungan latar belakang dinamis yang memiliki fluktuasi periodik (seperti permukaan air yang berombak), \textit{Gaussian Mixture Model} (GMM) memodelkan densitas probabilitas setiap nilai piksel sebagai kombinasi linear dari $K$ distribusi Gaussian:
\begin{equation}
    P(I_t) = \sum_{k=1}^K \omega_{k,t} \cdot \eta(I_t; \mu_{k,t}, \Sigma_{k,t})
    \label{eq:gmm}
\end{equation}
Meskipun GMM tangguh mengasimilasi riak air statis berulang, model ini memiliki inersia adaptasi yang relatif lambat terhadap event kemunculan sangat singkat (3--6 frame), sehingga rentan memicu fenomena \textit{ghosting} atau melewatkan puncak cipratan makan lele~\cite{Aris2023}.

\subsubsection{Deteksi Berbasis Deep Learning}
Dalam beberapa tahun terakhir, arsitektur \textit{Convolutional Neural Networks} (CNN) dan YOLO (\textit{You Only Look Once}) nano telah diadaptasi untuk mendeteksi tubuh ikan dan riak pakan secara langsung~\cite{Qian2025,Vijayalakshmi2025}. Model deep learning menawarkan kemampuan abstraksi fitur semantik yang unggul dalam membedakan objek asing (seperti daun hanyut atau bayangan burung). Namun, pengoperasian model deteksi objek pada resolusi video tinggi pada prosesor ARM tanpa akselerator tensor khusus menimbulkan latensi eksekusi yang mendekati atau melampaui batas batas waktu frame rate 10 FPS ($\ge 100$ ms per inferensi).

\subsection{Arsitektur Komputasi Tepi (Edge Computing) Menggunakan Raspberry Pi}
Komputasi tepi (\textit{edge computing}) memindahkan beban pemrosesan data analitik dari peladen terpusat (*cloud/server*) langsung ke simpul perangkat pemroses lokal di dekat sumber sensor. Raspberry Pi 4B/5 yang mengusung prosesor quad-core ARM Cortex-A72/A76 menjadi platform edge yang ideal untuk sistem akuakultur terdistribusi karena memiliki konsumsi daya rendah ($<15$ Watt), antarmuka kamera pita CSI (\textit{Camera Serial Interface}) berlatensi rendah, dan dukungan pustaka optimasi numerik OpenCV serta NumPy berbasis instruksi ARM NEON.

Pada sistem pemantau kolam lele, kamera dipasang pada posisi pandang atas (\textit{top-view}) tegak lurus terhadap permukaan kolam pada ketinggian $1,5 - 2,0$ meter. Alur data citra beroperasi pada laju tangkap 10 FPS, di mana setiap frame segera diproses dalam memori RAM tanpa penulisan ke media penyimpanan sekunder (SD card) guna mencegah \textit{I/O bottleneck}. Output yang dihasilkan berupa metrik numerik indeks deviasi gerak $\Delta M_t$ yang ditransmisikan melalui pesan JSON ke broker MQTT, menjamin efisiensi bandwidth transmisi jaringan nirkabel.

\section{TINJAUAN PUSTAKA}
\label{sec:tinjauan-pustaka}

\subsection{Strategi Penelusuran dan Seleksi Literatur}
Penelusuran literatur dalam penelitian ini dirancang secara sistematis untuk memetakan perkembangan terkini (\textit{state-of-the-art}) metode deteksi gerak visual pada ranah pemantauan perilaku makan akuakultur. Penelusuran dilakukan pada basis data bereputasi internasional dan nasional meliputi Consensus AI Academic Search Engine, IEEE Xplore, ScienceDirect (Elsevier), SpringerLink, Nature Portfolio, Frontiers in Marine Science, dan jurnal nasional terakreditasi SINTA.

Strategi penelusuran memanfaatkan formulasi query berbasis operator Boolean yang menggabungkan kata kunci domain akuakultur dengan metodologi visi komputer:
\begin{itemize}
    \item \texttt{("motion detection" OR "activity detection") AND ("fish pond" OR "aquaculture")}
    \item \texttt{("frame differencing" OR "optical flow" OR "background subtraction") AND ("edge" OR "embedded" OR "Raspberry Pi")}
    \item \texttt{("feeding behavior" OR "appetite assessment") AND ("water splash" OR "surface motion")}
    \item \texttt{("outdoor pond" OR "anti-disturbance" OR "shadow removal" OR "water wave noise")}
\end{itemize}

Berdasarkan kriteria seleksi yang berfokus pada publikasi 5 tahun terakhir (2022--2026) dengan penambahan artikel fondasional perintis, diperoleh 19 artikel ilmiah teks lengkap (\textit{full-text PDF}) yang diverifikasi secara ketat integritas dokumennya.

\subsection{Ulasan Mendalam Artikel Penelitian Kunci}
Dari 19 artikel yang terverifikasi, dilakukan penelaahan mendalam terhadap 8 artikel yang memiliki keterkaitan metodologis paling erat dengan kendala pemantauan kolam lele ber-aerator:

\subsubsection{Fukae et al. (2023) --- Dense Optical Flow untuk Klasifikasi Perilaku Kawanan}
Fukae et al.~\cite{Fukae2023} mengusulkan klasifikasi perilaku kawanan ikan menggunakan algoritma \textit{Gunnar Farnebäck dense optical flow}. Vektor kecepatan perpindahan piksel diekstraksi secara agregat dari seluruh luasan frame video, kemudian dikonversi menjadi histogram fitur gerak spasio-temporal dan diklasifikasikan menggunakan \textit{Support Vector Machine} (SVM). Pendekatan ini mencapai akurasi klasifikasi sebesar 88,4\% dalam membedakan kondisi makan dan non-makan. Keunggulan utamanya adalah kemampuan mengekstraksi dinamika kawanan tanpa perlu mendeteksi atau melacak individu ikan. Kendati demikian, komputasi aliran optik rapat pada setiap piksel membutuhkan waktu 42 ms pada kartu grafis GPU dan membengkak hingga $>350$ ms per frame pada prosesor ARM, menjadikannya tidak layak dieksekusi secara real-time pada Raspberry Pi tanpa akselerator GPU.

\subsubsection{Dong et al. (2024) --- Model Three-Stream Network untuk Intensitas Makan}
Dong et al.~\cite{Dong2024} mengembangkan arsitektur jaringan tiga jalur (\textit{three-stream network}) yang mengintegrasikan informasi spasial, perbedaan frame temporal ($|I_t - I_{t-1}|$), dan representasi domain frekuensi berbasis \textit{Fast Fourier Transform} (FFT) untuk mengklasifikasi intensitas makan ikan menjadi 4 tingkatan dengan akurasi 94,6\%. Temuan penting dari penelitian ini adalah bahwa aliran selisih frame temporal memberikan kontribusi fitur paling diskriminatif dalam menangkap lonjakan gerak mendadak. Namun, penggabungan model berbobot 48,2 MB dengan 14,6 GFLOPS menuntut infrastruktur server peladen yang tidak dapat dipenuhi oleh perangkat komputasi tepi berdaya rendah.

\subsubsection{Zhang et al. (2024) --- RUSNet: Segmentasi Ikan Berbasis Seleksi Adaptif Optical Flow}
Zhang et al.~\cite{Zhang2024} merancang arsitektur RUSNet yang mengintegrasikan mekanisme seleksi vektor aliran optik adaptif. Algoritma ini secara cerdas memisahkan komponen vektor gerak acak akibat arus fluida air dari vektor gerak terarah yang dihasilkan oleh perenang aktif, menghasilkan nilai mIoU segmentasi sebesar 86,2\% pada video perairan dinamis. Konsep pemisahan vektor derau fluida ini menjadi rujukan kunci untuk membedakan riak periodik aerator dari cipratan ikan lapar.

\subsubsection{Aris et al. (2023) --- Tantangan Background Subtraction dan GMM pada Latar Dinamis}
Aris et al.~\cite{Aris2023} menyajikan survei komprehensif mengenai tantangan pengurangan latar belakang menggunakan \textit{Gaussian Mixture Models} (GMM). Kajian ini menyoroti bahwa pada latar belakang yang memiliki pergerakan periodik cepat (seperti gelombang air atau gelembung), GMM cenderung mengalami kegagalan adaptasi karena laju pembaruan parameter model latar tidak cukup responsif untuk menangkap objek transien yang hanya muncul sesaat (3--6 frame), sehingga memicu fenomena penyerapan objek ke latar belakang (\textit{absorption}).

\subsubsection{Qian et al. (2025) --- PM-YOLOv8 Real-Time untuk Deteksi Makan Ikan}
Qian et al.~\cite{Qian2025} memodifikasi arsitektur YOLOv8n dengan menyematkan \textit{Parallelized Patch-Aware Attention} (PPA) dan MPDIoU loss untuk mendeteksi perilaku makan ikan bertumpuk di kolam air tawar. Pada dataset FeedFishDatas (7.400 citra), model mencapai presisi 87,7\%, mAP@50 83,5\%, dan latensi inferensi 8,1 ms pada GPU pengujian. Penulis mencatat bahwa implementasi model pada perangkat berdaya rendah (\textit{edge deployment}) melalui kuantisasi dan pemangkasan masih merupakan rencana pekerjaan lanjutan (\textit{future work}), serta metode ini mengasumsikan tubuh ikan terlihat di permukaan.

\subsubsection{Vijayalakshmi \& Sasithradevi (2025) --- AquaYOLO pada Kolam Luar Ruangan}
Vijayalakshmi dan Sasithradevi~\cite{Vijayalakshmi2025} mengusulkan AquaYOLO yang mengoptimalkan arsitektur YOLO melalui fusi fitur multi-skala dan head deteksi presisi untuk memantau kolam budidaya luar ruangan (\textit{aquaculture ponds}). Diuji pada dataset DePondFi (8.150 citra dari kolam di India Selatan), model mencapai presisi 0,889, recall 0,848, dan mAP@50 0,909, membuktikan kelayakan modifikasi YOLO untuk deteksi objek ikan pada lingkungan kolam terbuka.

\subsubsection{Marliza \& Aisuwarya (2021) --- Otomasi Pakan Lele Berbasis Mini-PC di Indonesia}
Marliza dan Aisuwarya~\cite{Marliza2021} mengimplementasikan sistem pakan lele otomatis menggunakan pengolahan citra skala abu-abu dan \textit{background subtraction} pada Mini-PC. Penelitian ini membuktikan kelayakan implementasi perangkat berbiaya murah untuk peternak lele lokal. Namun, pengujian dilakukan pada akuarium skala laboratorium tanpa aerator dan tanpa fluktuasi cuaca luar ruangan, sehingga algoritma tersebut akan menghasilkan keputusan pakan terus-menerus jika dioperasikan pada kolam ber-aerator.

\subsubsection{Liu et al. (2014) \& Zhao et al. (2016) --- CVFAI dan Pemodelan Energi Kinetik}
Liu et al.~\cite{Liu2014} memperkenalkan indeks aktivitas makan \textit{Computer Vision-based Feeding Activity Index} (CVFAI) berbasis rasio piksel biner hasil perbedaan citra temporal terhadap total luasan kolam. Formulasi ini hanya membutuhkan waktu komputasi $<5$ ms pada CPU Raspberry Pi. Zhao et al.~\cite{Zhao2016} memperluas konsep tersebut menjadi pemodelan energi kinetik statistik. Indeks matematis ini sangat ideal sebagai metrik numerik telemetri edge, namun memerlukan modifikasi kalibrasi untuk mengurangkan nilai ambang dasar (\textit{baseline}) yang dihasilkan oleh aerator.

\subsection{Sintesis Tematik Literatur}
Berdasarkan penelaahan terhadap keseluruhan 19 artikel, literatur yang dikaji terkonvergensi ke dalam tiga tema besar:
\begin{enumerate}
    \item \textbf{Tema 1 --- Deteksi Perilaku Makan via Deep Learning}: Mengandalkan model deteksi objek atau ekstraksi fitur berbasis konvolusi/Transformer (Qian et al.~\cite{Qian2025}, Hu et al.~\cite{Hu2026}, Zhao et al.~\cite{Zhao2025}, Dong et al.~\cite{Dong2024}, Vijayalakshmi et al.~\cite{Vijayalakshmi2025}). Pendekatan ini memberikan akurasi tinggi namun memiliki beban komputasi besar dan rentan terhadap keterbatasan visibilitas air keruh.
    \item \textbf{Tema 2 --- Kuantifikasi Gerakan dan Dinamika Temporal}: Mengukur kecepatan dan energi kinetik renang ikan secara non-segmentasi menggunakan optical flow dan selisih frame temporal (Fukae et al.~\cite{Fukae2023}, Alahmad et al.~\cite{Alahmad2025}, Zhao et al.~\cite{Zhao2016}, Zhang et al.~\cite{Zhang2024}). Pendekatan ini sangat efektif menangkap dinamika kelompok ikan secara agregat.
    \item \textbf{Tema 3 --- Pemodelan Latar Belakang dan Otomasi Berbasis Edge}: Menitikberatkan pada pengurangan latar belakang adaptif, pemodelan derau dinamis, dan perancangan sistem tertanam berbiaya rendah (Aris et al.~\cite{Aris2023}, Marliza et al.~\cite{Marliza2021}, Atoum et al.~\cite{Atoum2015}, Zhou et al.~\cite{Zhou2018}, Liu et al.~\cite{Liu2014}, Zhao et al.~\cite{Zhao2017}). Kelompok ini menjadi pilar utama perancangan algoritma ringan yang mampu mengekstraksi metrik numerik pada perangkat tepi.
\end{enumerate}

\section{PERBANDINGAN PENELITIAN}
\label{sec:perbandingan}

% Bab ini memuat analisis perbandingan sistematis antar-artikel ilmiah yang diteliti.
Perbandingan sistematis terhadap penelitian-penelitian terdahulu disajikan pada Tabel~\ref{tab:article-comparison}.

% Article Comparison Table (Longtable for multi-page support)
\begin{small}
\begin{longtable}{p{2.6cm} p{3.2cm} p{2.8cm} p{3.4cm} p{2.6cm}}
\caption{Matriks Perbandingan Artikel Penelitian Utama Hasil Sintesis Consensus\label{tab:article-comparison}} \\
\toprule
\textbf{Artikel} & \textbf{Metode / Pendekatan} & \textbf{Dataset / Lingkungan} & \textbf{Hasil Utama} & \textbf{Keterbatasan} \\
\midrule
\endfirsthead

\caption[]{Matriks Perbandingan Artikel Penelitian Utama (\textit{Lanjutan})} \\
\toprule
\textbf{Artikel} & \textbf{Metode / Pendekatan} & \textbf{Dataset / Lingkungan} & \textbf{Hasil Utama} & \textbf{Keterbatasan} \\
\midrule
\endhead

\midrule
\multicolumn{5}{r}{\textit{Bersambung ke halaman berikutnya\dots}} \\
\bottomrule
\endfoot

\bottomrule
\endlastfoot

{Qian et al.~\cite{Qian2025}} & PM-YOLO (YOLOv8n + Patch Attention) & Kolam air tawar (FeedFishDatas) & Presisi 87,7\%, mAP 83,5\%, latensi 8,1 ms & Fokus tubuh ikan di permukaan; rentan air sangat keruh \\
\addlinespace

{Hu et al.~\cite{Hu2026}} & HVIT (CNN + ViT + LSTM) & Tangki RAS industri (Largemouth Bass) & Akurasi $>$98\% kuantifikasi intensitas makan & Diuji pada RAS terkontrol; komputasi berat untuk edge \\
\addlinespace

{Alahmad et al.~\cite{Alahmad2025}} & YOLOv8 + Estimasi Kecepatan Renang & Video akuakultur laut & Akurasi 85\%; beda kecepatan renang lapar vs kenyang & Mengasumsikan tubuh ikan terlihat; belum uji riak keruh \\
\addlinespace

{Fukae et al.~\cite{Fukae2023}} & Dense Optical Flow + Histogram + SVM & Kolam budidaya ikan & Klasifikasi makan vs tidak makan secara agregat & Komputasi optical flow padat; belum atasi gelembung aerator \\
\addlinespace

{Zhao et al.~\cite{Zhao2025}} & Two-Stage Net (Pose + GCN) & Tangki akuakultur industri & Akurasi kenyang 98,1\%, reduksi parameter 31,4\% & Bergantung pada pose fisik ikan yang terlihat jelas \\
\addlinespace

{Li et al.~\cite{Li2020}} & Tinjauan Sistematis Pengenalan Makan & Sistem budidaya global & Peta komparatif visual, akustik, dan biosensor & Tinjauan literatur tanpa implementasi sistem outdoor \\
\addlinespace

{Marliza et al.~\cite{Marliza2021}} & Pengolahan Citra Digital + Background Sub & Akuarium / kolam ikan air tawar & Pemberian pakan otomatis saat ikan aktif di permukaan & Skala laboratorium; belum tangani aerator dan cuaca outdoor \\
\addlinespace

{Atoum et al.~\cite{Atoum2015}} & Analisis Gerak Kuantitatif Temporal & Tangki budidaya densitas tinggi & Otomasi kontrol pakan adaptif mengurangi limbah & Diuji pada pencahayaan konstan tanpa turbulensi aerator \\
\addlinespace

{Al-Abri et al.~\cite{Al-Abri2025}} & Tinjauan Computer Vision Akuakultur & Literatur global AI dan IoT & Pemetaan tren metode deteksi dan klasifikasi pakan & Tinjauan survei tanpa pengujian empiris khusus lele \\
\addlinespace

{Zhou et al.~\cite{Zhou2018}} & Near-Infrared Vision + ANFIS Neuro-Fuzzy & Sistem sirkulasi akuakultur & Turunkan FCR 10,77\% dan polusi air 9,8\% & Gunakan kamera NIR khusus; belum tangani aerator outdoor \\
\addlinespace

{Liu et al.~\cite{Liu2014}} & Feeding Activity Index (CVFAI) & Tangki RAS salmon Atlantik & Deteksi aktivitas non-kontak berbiaya rendah & Diuji pada air jernih; belum teruji pada air kolam keruh \\
\addlinespace

{Zhao et al.~\cite{Zhao2017}} & Analisis Perilaku Kolektif Spontan & Tangki sirkulasi akuakultur & Mampu menilai nafsu makan sebelum pakan ditebar & Asumsi latar belakang statis tanpa aerator gelembung kontinu \\
\addlinespace

{Zhao et al.~\cite{Zhao2016}} & Pemodelan Energi Kinetik Spasial-Temporal & Tangki RAS ikan nila & Deteksi perilaku abnormal dan perubahan aktivitas & Belum mempertimbangkan turbulensi konstan aerator outdoor \\
\addlinespace

{Dong et al.~\cite{Dong2024}} & Three-Stream Network (Spatial, Temporal, Frequency) & Akuakultur industri (FeedFishNet) & Klasifikasi intensitas makan presisi tinggi multi-skala & Membutuhkan akselerasi GPU; belum diuji pada single-board computer \\
\addlinespace

{Zhang et al.~\cite{Zhang2024}} & RUSNet (Adaptive Optical Flow Selection) & Video lingkungan akuatik dinamis & Segmentasi objek ikan adaptif terhadap degradasi visual & Fokus pada segmentasi under-water, bukan riak permukaan \\
\addlinespace

{Aris et al.~\cite{Aris2023}} & Survei Background Subtraction \& GMM & Dataset video dinamis multikondisi & Taksonomi komparatif adaptasi latar belakang dinamis & Kajian teoritis; butuh penyesuaian khusus gelombang aerator \\
\addlinespace

{Kamal et al.~\cite{Kamal2025}} & Akustik Pasif (Hydrophone) + Machine Learning & Kolam lele intensif (\textit{Clarias gariepinus}) & Otomasi pakan lele densitas tinggi berbasis suara kunyahan & Menggunakan sensor audio/suara, bukan visi kamera komputer \\
\addlinespace

{Yang et al.~\cite{Yang2025}} & Tinjauan Komprehensif State-of-the-Art 2025 & Domain computer vision akuakultur & Analisis kelemahan model deep learning terkini & Tinjauan survei sekunder tanpa data pengujian empiris langsung \\
\addlinespace

{Vijayalakshmi et al.~\cite{Vijayalakshmi2025}} & AquaYOLO (Backbone Teroptimasi \& Attention) & Kolam luar ruangan (\textit{outdoor ponds}) & Deteksi ikan akurat pada lingkungan terbuka dengan derau visual & Belum memodelkan dinamika temporal riak makan sub-detik \\

\end{longtable}
\end{small}


[Uraikan analisis dan pembahasan mengenai matriks perbandingan penelitian di atas...]

\section{KESENJANGAN PENELITIAN (\textit{RESEARCH GAP})}
\label{sec:research-gap}

% Identifikasi celah atau kesenjangan yang belum terselesaikan oleh penelitian terdahulu.
Berdasarkan analisis kritis terhadap literatur yang ditinjau, ditemukan beberapa kesenjangan penelitian (\textit{research gap}) utama:

\subsection{Kesenjangan Metodologis dan Integrasi}
[Uraikan kesenjangan metodologis, keterbatasan teknis, atau kekurangan integrasi dari penelitian terdahulu...]

\subsection{Kesenjangan Evaluasi dan Pengujian}
[Uraikan kekurangan atau celah dalam hal evaluasi empiris, dataset, atau pengujian performa...]

\section{POSISI PENELITIAN (\textit{THESIS POSITIONING})}
\label{sec:positioning}

% Jelaskan posisi penelitian Anda relatif terhadap penelitian-penelitian sebelumnya.
\subsection{Pemosisian Terhadap Penelitian Terkait}
[Jelaskan posisi penelitian ini di antara penelitian-penelitian yang sudah ada sebelumnya...]

\subsection{Kontribusi dan Kebaruan Penelitian}
Penelitian ini diharapkan memberikan kontribusi utama sebagai berikut:
\begin{enumerate}
    \item [Kontribusi atau kebaruan 1]
    \item [Kontribusi atau kebaruan 2]
\end{enumerate}

\section{KESIMPULAN DAN REKOMENDASI METODE}
\label{sec:kesimpulan}

\subsection{Kesimpulan}
Studi literatur yang dilakukan terhadap 19 artikel ilmiah terverifikasi (dengan lebih dari 63\% terbit pada rentang 2023--2026) menghasilkan sejumlah kesimpulan penting:
\begin{enumerate}
    \item Pemantauan perilaku makan ikan secara non-kontak berbasis visi komputer telah terbukti efektif dalam literatur global untuk menekan rasio konversi pakan (FCR), meningkatkan kualitas air, dan mencegah kanibalisme lele.
    \item Terdapat kesenjangan penelitian mendasar (\textit{continuous aeration baseline gap}) di mana sebagian besar metode terdahulu mengasumsikan permukaan air kolam yang tenang. Pada kolam budidaya lele luar ruangan yang dilengkapi aerator mekanis 24 jam, metode deteksi gerak konvensional mengalami kegagalan akibat positif palsu yang dipicu oleh turbulensi gelembung air.
    \item Karakteristik kemunculan makan lele pada air keruh bersifat sangat singkat dan impulsif (berlangsung hanya 3 hingga 6 frame pada kecepatan rekam 10 FPS atau durasi 0,3--0,6 detik). Model deteksi deep learning berbasis bentuk tubuh dan dense optical flow penuh terbukti terlalu berat dan lambat untuk dioperasikan pada perangkat edge Raspberry Pi tanpa kartu grafis.
    \item Diperlukan metode hibrida yang mengombinasikan efisiensi komputasi perbedaan frame temporal ($O(N)$) dengan pemodelan statistik varians latar belakang untuk mereduksi derau aerator secara dinamis.
\end{enumerate}

\subsection{Rekomendasi Metode untuk Diuji Coba}
Berdasarkan sintesis komparatif yang telah dilakukan, dirumuskan dua rekomendasi metode untuk diimplementasikan dan diuji coba pada tahap eksperimen skripsi:

\subsubsection{Metode Utama (Rekomendasi 1): Adaptive Temporal Variance Subtraction}
Metode utama yang direkomendasikan adalah pendekatan hibrida dua tahap:
\begin{enumerate}
    \item \textbf{Tahap Kalibrasi Baseline}: Menghitung peta varians temporal lokal $\sigma_{\text{aerator}}^2(x,y)$ dari sekuens air kolam saat aerator beroperasi untuk memetakan sebaran turbulensi gelembung air.
    \item \textbf{Tahap Deteksi Transien}: Menerapkan operasi perbedaan citra temporal absolut $|I_t - I_{t-1}|$ dengan ambang batas adaptif berbasis peta varians $\Theta(x,y) = \gamma \cdot \sigma_{\text{aerator}}(x,y) + \tau_0$.
    \item \textbf{Ekstraksi Metrik Numerik}: Menghitung indeks deviasi aktivitas $\Delta M_t$ pada CPU Raspberry Pi 4B/5 yang ditargetkan mampu memproses citra pada laju tangkap 10 FPS, lalu mentransmisikannya sebagai payload data numerik via protokol MQTT.
\end{enumerate}
Pendekatan ini menjadi alternatif prioritas untuk dieksplorasi karena memiliki kompleksitas komputasi yang relatif ringan, sensitif terhadap perubahan mendadak, dan tidak menuntut dataset beranotasi besar.

\subsubsection{Arah Alternatif untuk Dicoba: Quantized Lightweight YOLO + FSM}
Sebagai model pembanding alternatif yang dapat dieksplorasi jika sumber daya komputasi memungkinkan:
\begin{enumerate}
    \item Menguji model deteksi objek ringan (YOLOv8-nano) yang dikuantisasi ke format INT8 (NCNN atau TensorFlow Lite) untuk mengevaluasi kemampuan deteksi cipratan air (\textit{feeding splash}) atau bagian tubuh ikan di permukaan.
    \item Memvalidasi kontinuitas deteksi menggunakan mesin status terhingga (\textit{Finite State Machine}) temporal untuk mengurangi kesalahan deteksi akibat gangguan acak.
\end{enumerate}

\subsection{Saran dan Arah Pengembangan}
Pada tahap pengujian eksperimental berikutnya, disarankan untuk:
\begin{itemize}
    \item Menguji ketahanan algoritma terhadap perubahan pencahayaan ekstrem (kondisi terik matahari siang hari vs. pencahayaan lampu buatan malam hari).
    \item Melakukan kalibrasi nilai sensitivitas ambang batas $\gamma$ terhadap variasi debit aerasi dan kepadatan populasi ikan lele di kolam uji.
\end{itemize}


\newpage

% ==================== DAFTAR PUSTAKA ====================
\addcontentsline{toc}{section}{DAFTAR PUSTAKA}
\nocite{*}
\bibliographystyle{plainnat}
\bibliography{references}

\end{document}
